\section{Syntax at the Sentence Level}
\label{sec:sentence-level-syntax}
\slidesrc{Syntax: sentence, rules: 1--17; Syntax: sentence, HMM: 1--20}

Once a text is split into words, the next questions concern whole sentences: where does a
sentence end, what grammatical category does each word have in context, and how do the
words group into a structure? This section covers the sentence-level tasks of the
syntactic level (Section~\ref{sec:course-map}):
\begin{enumerate}[nosep]
  \item \textbf{sentence breaking} (Section~\ref{sec:sentence-breaking});
  \item \textbf{part-of-speech tagging} and the Penn Treebank tagset
        (Section~\ref{sec:pos-tagging});
  \item \textbf{grammars and parsing}: Chomsky's generative grammar, the Chomsky
        hierarchy, a small grammar for English, Definite Clause Grammars
        (Section~\ref{sec:grammars-parsing});
  \item \textbf{statistical POS tagging} with Hidden Markov Models and the Viterbi
        algorithm (Section~\ref{sec:hmm-pos-tagging}).
\end{enumerate}
The first three are \emph{rule-based} (deductive) approaches, the last one is
\emph{statistical} (inductive): two eras of NLP (Section~\ref{sec:nlp-history}) side by
side.

\subsection{Sentence Breaking}
\label{sec:sentence-breaking}
\slidesrc{Syntax: sentence, rules: 2--4}

\begin{nbdef}{Sentence boundary disambiguation}
\textbf{Sentence boundary disambiguation} (SBD), also called \emph{sentence breaking}
or \emph{sentence boundary detection}, is the problem of deciding where sentences begin
and end.
\end{nbdef}

Many NLP tools need their input divided into sentences, but punctuation is ambiguous:
\begin{itemize}[nosep]
  \item a \textbf{period} may mark an abbreviation (\emph{Mr.}, \emph{etc.}), a decimal
        point (\emph{3.14}), an ellipsis (\emph{\ldots}) or be part of an e-mail
        address, rather than the end of a sentence. About \textbf{47\%} of the periods
        in the Wall Street Journal corpus mark abbreviations;
  \item \textbf{question and exclamation marks} appear in embedded quotations, emoticons,
        computer code and slang.
\end{itemize}
Languages like Japanese and Chinese have unambiguous markers for the end of a sentence.

\begin{nbfact}{The standard rule-based approach}
To decide whether a token ends a sentence:
\begin{enumerate}[label=(\alph*), nosep]
  \item if it is a period, it ends a sentence;
  \item but if the preceding token is in a hand-compiled list of
        \textbf{abbreviations}, it does not;
  \item and if the next token is \textbf{capitalised}, then it does end a sentence.
\end{enumerate}
This simple strategy gets about \textbf{95\%} of sentences right.
\end{nbfact}

\begin{nbexample}{where the rules break}
\emph{Mr. Smith paid 3.14 dollars. He left at 5 p.m. Then he slept.}
\begin{itemize}[nosep]
  \item \emph{Mr.}: abbreviation, rule (b): not a boundary. Correct.
  \item \emph{3.14}: the period is inside a token, not a token of its own. Correct.
  \item \emph{p.m. Then}: \emph{p.m.} is an abbreviation (b says ``no boundary''), but
        the next word is capitalised (c says ``boundary''). Here the sentence
        \emph{does} end: the period of an abbreviation can also be the final period.
\end{itemize}
NLTK's pre-trained sentence splitter returns \emph{[``Mr. Smith paid 3.14 dollars.'',
``He left at 5 p.m. Then he slept.'']}: it gets the first two cases right and misses
exactly the third. The remaining 5\% of errors are cases like this one, which need more
context (or a statistical model) to resolve.
\end{nbexample}

\subsection{Part-of-Speech Tagging}
\label{sec:pos-tagging}
\slidesrc{Syntax: sentence, rules: 5--6}

\begin{nbdef}{Part-of-speech tagging}
\textbf{Part-of-speech tagging} (POS tagging, POST), also called \emph{grammatical
tagging} or \emph{word-category disambiguation}, is the process of marking each word of
a text with its part of speech, based on both its \textbf{definition} and its
\textbf{context}, i.e.\ its relationship with adjacent and related words in a phrase,
sentence or paragraph.
\end{nbdef}

(If the categories themselves are fuzzy, see the recap with the Italian equivalents in
Section~\ref{sec:parts-of-speech}.) The word ``context'' is the whole point: many words have several possible tags, and only
the sentence tells which one is right.

\begin{nbexample}{the same word, different tags}
\begin{itemize}[nosep]
  \item \emph{I \underline{book}/VB a flight} \quad vs.\quad \emph{I read a
        \underline{book}/NN}.
  \item \emph{\underline{Time}/NN \underline{flies}/VBZ \underline{like}/IN an arrow}
        (time passes quickly) \quad vs.\quad \emph{\underline{Time}/VB
        \underline{flies}/NNS \underline{like}/IN an arrow} (measure the speed of flies as
        you would an arrow) \quad vs.\quad \emph{Time/NN \underline{flies}/NNS
        \underline{like}/VBP an arrow} (``time flies'', a kind of fly, are fond of an
        arrow).
\end{itemize}
\end{nbexample}

\subsubsection{The Penn Treebank tagset}
\slidesrc{Syntax: sentence, rules: 6}

The most widely used set of tags for English is the \textbf{Penn Treebank} tagset, with
36 word tags:

\begin{center}
\footnotesize
\begin{tabular}{@{}r l l @{\hspace{1.2cm}} r l l@{}}
\toprule
 & \textbf{Tag} & \textbf{Description} & & \textbf{Tag} & \textbf{Description} \\
\midrule
1  & CC   & coordinating conjunction       & 19 & PRP\$ & possessive pronoun \\
2  & CD   & cardinal number                & 20 & RB    & adverb \\
3  & DT   & determiner                     & 21 & RBR   & adverb, comparative \\
4  & EX   & existential \emph{there}       & 22 & RBS   & adverb, superlative \\
5  & FW   & foreign word                   & 23 & RP    & particle \\
6  & IN   & preposition / subord.\ conj.   & 24 & SYM   & symbol \\
7  & JJ   & adjective                      & 25 & TO    & \emph{to} \\
8  & JJR  & adjective, comparative         & 26 & UH    & interjection \\
9  & JJS  & adjective, superlative         & 27 & VB    & verb, base form \\
10 & LS   & list item marker               & 28 & VBD   & verb, past tense \\
11 & MD   & modal                          & 29 & VBG   & verb, gerund / present participle \\
12 & NN   & noun, singular or mass         & 30 & VBN   & verb, past participle \\
13 & NNS  & noun, plural                   & 31 & VBP   & verb, non-3rd person sing.\ present \\
14 & NNP  & proper noun, singular          & 32 & VBZ   & verb, 3rd person sing.\ present \\
15 & NNPS & proper noun, plural            & 33 & WDT   & \emph{wh}-determiner \\
16 & PDT  & predeterminer                  & 34 & WP    & \emph{wh}-pronoun \\
17 & POS  & possessive ending              & 35 & WP\$  & possessive \emph{wh}-pronoun \\
18 & PRP  & personal pronoun               & 36 & WRB   & \emph{wh}-adverb \\
\bottomrule
\end{tabular}
\captionof{table}{The Penn Treebank POS tagset.}
\end{center}

\begin{nbintuition}[reading the tags]
The tags encode morphology (Section~\ref{sec:morphology}): NN/NNS differ by the plural
inflection, VB/VBD/VBG/VBN/VBP/VBZ by the verb's inflection, JJ/JJR/JJS by comparison.
So a tagger is also, partly, a morphological analyser, and the tag is what a lemmatiser
needs to choose the right lemma (Section~\ref{sec:stemming-lemmatization}).
\end{nbintuition}

\begin{nbexample}{a tagged sentence}
\emph{The/DT boy/NN 's/POS cars/NNS are/VBP different/JJ colors/NNS ./.}
\end{nbexample}

Online taggers to try: \url{http://parts-of-speech.info/} and Stanford CoreNLP,
\url{https://corenlp.run/}.

There are two families of approaches, both covered in this course:
\begin{itemize}[nosep]
  \item \textbf{rule-based} (deductive): a grammar and a parser decide the tags as a
        by-product of the structure of the sentence (Section~\ref{sec:grammars-parsing});
  \item \textbf{statistical} (inductive): a model learnt from a tagged corpus picks the
        most probable sequence of tags (Section~\ref{sec:hmm-pos-tagging}).
\end{itemize}

\subsection{Grammars and Parsing}
\label{sec:grammars-parsing}
\slidesrc{Syntax: sentence, rules: 7--17}

\begin{nbdef}{Parsing}
In computational linguistics, \textbf{parsing} is the formal analysis by a computer of a
sentence (or other string of words) into its \textbf{constituents}, resulting in a
\textbf{parse tree} that shows their syntactic relations to each other. The tree may also
contain semantic and other information.
\end{nbdef}

Parsing is one of the possible ways to address the POS problem: once the sentence has a
tree, every word has a category (the node right above it).

\subsubsection{Chomsky and generative grammar}
\slidesrc{Syntax: sentence, rules: 8--9}

In \emph{Syntactic Structures} (1957), Noam Chomsky proposed the theory of
\textbf{generative grammar}: a finite set of rules that recursively ``specify'' or
``generate'' all the well-formed expressions of a natural language (and only them). It
is considered one of the most significant contributions to theoretical linguistics of
the 20th century.

Recursion is what links generative grammar to the question of
Section~\ref{sec:human-vs-animal}. ``Animal languages'' are forms of animal communication
that show similarities to human language: the inventory of signs may be large, the signs
relatively arbitrary, and they seem to be produced with some volition. But researchers
agree that they are not as complex or expressive as human language: they lack the
\emph{creation of new patterns of signs under varied circumstances}, i.e.\ the
productivity that recursive rules provide.

\subsubsection{Grammars and the Chomsky hierarchy}
\slidesrc{Syntax: sentence, rules: 10--13}

\begin{nbdef}{Grammar and formal language}
A \textbf{grammar} specifies the compositional structure of complex messages. A
\textbf{formal language} is a set of strings of \textbf{terminal symbols} (the words).
Each string of the language can be analysed (recognised) or generated by the grammar,
which is a set of \textbf{rewrite rules}, e.g.
\begin{align*}
  S &\to NP\ VP\\
  Article &\to \textbf{the} \mid \textbf{a} \mid \textbf{an} \mid \dots
\end{align*}
Here $S$ is the \textbf{sentence} (start) symbol; $NP$ (noun phrase), $VP$ (verb
phrase) and $Article$ are \textbf{nonterminals}; \textbf{the}, \textbf{a}, \textbf{an}
are terminals; the bar $\mid$ separates alternatives.
\end{nbdef}

Restricting the \emph{form} of the rules gives classes of grammars of increasing power,
the \textbf{Chomsky--Schützenberger hierarchy}:

\begin{center}
\small
\begin{tabular}{@{}l l l l@{}}
\toprule
\textbf{Grammar} & \textbf{Languages} & \textbf{Automaton} & \textbf{Rules} \\
\midrule
Type-0 & recursively enumerable & Turing machine & no restrictions \\
Type-1 & context-sensitive & linear-bounded automaton & $\alpha A \beta \to \alpha\gamma\beta$ \\
Type-2 & context-free & (non-det.) pushdown automaton & $A \to \gamma$ \\
Type-3 & regular & finite-state automaton & $A \to aB$, \ $A \to a$ \\
\bottomrule
\end{tabular}
\captionof{table}{The Chomsky hierarchy. $A, B$: nonterminals; $a$: a terminal;
$\alpha, \beta, \gamma$: strings of terminals and nonterminals.}
\end{center}

\begin{center}
\begin{tikzpicture}
  \draw[fill=nbblue!6,   draw=nbblue]   (0,0.25) ellipse (4.4 and 2.3);
  \draw[fill=nbgreen!8,  draw=nbgreen]  (0,-0.35) ellipse (3.3 and 1.55);
  \draw[fill=nborange!10,draw=nborange] (0,-0.7) ellipse (2.4 and 1.1);
  \draw[fill=nbpurple!10,draw=nbpurple] (0,-1.0) ellipse (1.4 and 0.65);
  \node[font=\small] at (0,1.85) {recursively enumerable};
  \node[font=\small] at (0,0.8) {context-sensitive};
  \node[font=\small] at (0,0.0) {context-free};
  \node[font=\small] at (0,-1.0) {regular};
\end{tikzpicture}
\captionof{figure}{Each class contains the ones below it. Regular expressions
(Section~\ref{sec:regular-expressions}) describe exactly the regular languages.}
\end{center}

In words:
\begin{itemize}[nosep]
  \item \textbf{Regular}: \emph{nonterminal} $\to$ \emph{terminal} [\emph{nonterminal}],
        e.g.\ $S \to \mathbf{a}S$, $S \to \Lambda$ (the empty string): any number of
        \textbf{a}. Recognised with a fixed, finite memory.
  \item \textbf{Context-free}: \emph{nonterminal} $\to$ anything, e.g.\
        $S \to \mathbf{a}S\mathbf{b}$: generates $\mathbf{a}^n\mathbf{b}^n$, which needs
        to \emph{count} the a's, impossible with finite memory (a stack suffices).
  \item \textbf{Context-sensitive}: more nonterminals on the left-hand side, e.g.\
        $ASB \to AA\mathbf{a}BB$: $S$ can be rewritten only when surrounded by $A$ and
        $B$.
  \item \textbf{Recursively enumerable}: no constraints.
\end{itemize}

\begin{nbwarning}[``natural languages are probably context-free'']
The slides say that natural languages are probably context-free, hence parsable in real
time. This is a good working assumption: context-free grammars cover most syntactic
constructions and are the basis of most parsers. Strictly speaking, however, a few
constructions are known to go beyond context-free power, e.g.\ the cross-serial
dependencies of Swiss German (Shieber, 1985). Linguists therefore say that natural
languages are \emph{mildly} context-sensitive.
\end{nbwarning}

\subsubsection{A grammar for the Wumpus world}
\slidesrc{Syntax: sentence, rules: 14--15}

The Wumpus world is a toy game from AI textbooks: an agent explores a cave with pits and
a monster, the wumpus, and perceives stenches and breezes. A tiny fragment of English for
talking about it consists of a \textbf{lexicon} and a \textbf{grammar}.

\begin{center}
\small
\begin{tabular}{@{}r c l@{}}
\multicolumn{3}{@{}l}{\textbf{Lexicon}}\\
\midrule
\emph{Noun}        & $\to$ & stench $\mid$ breeze $\mid$ glitter $\mid$ nothing $\mid$ wumpus $\mid$ pit $\mid$ pits $\mid$ gold $\mid$ east $\mid \dots$ \\
\emph{Verb}        & $\to$ & is $\mid$ see $\mid$ smell $\mid$ shoot $\mid$ feel $\mid$ stinks $\mid$ go $\mid$ grab $\mid$ carry $\mid$ kill $\mid$ turn $\mid \dots$ \\
\emph{Adjective}   & $\to$ & right $\mid$ left $\mid$ east $\mid$ south $\mid$ back $\mid$ smelly $\mid \dots$ \\
\emph{Adverb}      & $\to$ & here $\mid$ there $\mid$ nearby $\mid$ ahead $\mid$ right $\mid$ left $\mid$ east $\mid \dots$ \\
\emph{Pronoun}     & $\to$ & me $\mid$ you $\mid$ I $\mid$ it $\mid \dots$ \\
\emph{Name}        & $\to$ & John $\mid$ Mary $\mid$ Boston $\mid \dots$ \\
\emph{Article}     & $\to$ & the $\mid$ a $\mid$ an $\mid \dots$ \\
\emph{Preposition} & $\to$ & to $\mid$ in $\mid$ on $\mid$ near $\mid \dots$ \\
\emph{RelPronoun}  & $\to$ & that $\mid$ which $\mid \dots$ \\
\end{tabular}
\end{center}

The lexicon is divided into \textbf{closed classes} (articles, pronouns, prepositions:
a small fixed set, new members are almost never added) and \textbf{open classes} (nouns,
verbs, adjectives, adverbs: new words appear all the time).

\begin{center}
\small
\begin{tabular}{@{}r c l l@{}}
\multicolumn{4}{@{}l}{\textbf{Grammar}}\\
\midrule
$S$         & $\to$ & $NP\ VP$                & I + feel a breeze \\[3pt]
$Subject$   & $\to$ & $Pronoun$               & I \\
            & $\mid$ & $Noun$                 & pits \\
            & $\mid$ & $Article\ Noun$        & the + wumpus \\
            & $\mid$ & $Article\ Adjective\ Noun$ & the + smelly + wumpus \\[3pt]
$NP$        & $\to$ & $Subject$               & the wumpus \\
            & $\mid$ & $Subject\ PP$          & the wumpus + to the east \\
            & $\mid$ & $Subject\ RelClause$   & the wumpus + that is smelly \\[3pt]
$VP$        & $\to$ & $Verb$                  & stinks \\
            & $\mid$ & $Verb\ NP$             & feel + a breeze \\
            & $\mid$ & $Verb\ Adjective$      & is + smelly \\
            & $\mid$ & $Verb\ PP$             & turn + to the east \\
            & $\mid$ & $Verb\ Adverb$         & go + ahead \\[3pt]
$PP$        & $\to$ & $Preposition\ NP$       & to + the east \\
$RelClause$ & $\to$ & $RelPronoun\ VP$        & that + is smelly \\
\end{tabular}
\end{center}

All rules have a single nonterminal on the left: this is a \textbf{context-free}
grammar. Note the recursion: $NP \to Subject\ PP$ and $PP \to Preposition\ NP$, so noun
phrases can nest without limit (\emph{the wumpus to the east of the pit near the gold
\ldots}).

\begin{nbexample}{parsing ``I feel a breeze''}
\begin{center}
\begin{tikzpicture}[level distance=0.95cm,
    level 1/.style={sibling distance=4.2cm},
    level 2/.style={sibling distance=2.2cm},
    level 3/.style={sibling distance=1.4cm},
    every node/.style={font=\small}]
  \node {$S$}
    child {node {$NP$}
      child {node {$Subject$}
        child {node {$Pronoun$}
          child {node[nbblue] {I}}}}}
    child {node {$VP$}
      child {node {$Verb$}
        child {node[nbblue] {feel}}}
      child {node {$NP$}
        child {node {$Subject$}
          child {node {$Article$} child {node[nbblue] {a}}}
          child {node {$Noun$} child {node[nbblue] {breeze}}}}}};
\end{tikzpicture}
\captionof{figure}{Parse tree. The node right above each word is its part of speech:
\emph{I}/Pronoun, \emph{feel}/Verb, \emph{a}/Article, \emph{breeze}/Noun.}
\end{center}
Derivation, always rewriting the leftmost nonterminal:
$S \Rightarrow NP\ VP \Rightarrow Subject\ VP \Rightarrow Pronoun\ VP \Rightarrow
\text{I}\ VP \Rightarrow \text{I}\ Verb\ NP \Rightarrow \text{I feel}\ NP \Rightarrow
\text{I feel}\ Subject \Rightarrow \text{I feel}\ Article\ Noun \Rightarrow
\text{I feel a breeze}$.
\end{nbexample}

\begin{nbwarning}[this grammar overgenerates]
The same rules also generate \emph{*I stinks}, \emph{*me feel a breeze} and \emph{*the
wumpus feel a breeze}: nothing checks that subject and verb agree in person and number,
or that a pronoun in subject position is \emph{I} and not \emph{me}. Adding a separate
nonterminal for every combination (singular/plural, subject/object, \ldots) would make the
grammar explode. Definite Clause Grammars solve this with \emph{arguments}.
\end{nbwarning}

\subsubsection{Definite Clause Grammars}
\slidesrc{Syntax: sentence, rules: 16--17}

\begin{nbdef}{Definite Clause Grammars (DCG)}
\textbf{DCGs} are a notation for grammars natively integrated in the Prolog programming
language.
\begin{itemize}[nosep]
  \item They represent context-free grammars\ldots
  \item \ldots but nonterminals can have \textbf{arguments} (compare with \emph{attribute
        grammars}), which can express context-dependent constraints: this makes the
        formalism extremely powerful.
  \item Being Prolog, the same grammar works both as a \textbf{recogniser} (is this
        sentence grammatical?) and as a \textbf{generator} (list all grammatical
        sentences), and arbitrary reasoning can be added on top of it (``but, to your
        relief, we will not!'').
\end{itemize}
\end{nbdef}

\begin{nbexample}{agreement with an argument}
\begin{minted}{prolog}
s          --> np(Num), vp(Num).          % subject and verb share Num
np(Num)    --> article, noun(Num).
vp(Num)    --> verb(Num).
article    --> [the].
noun(sing) --> [wumpus].     noun(plur) --> [wumpuses].
verb(sing) --> [stinks].     verb(plur) --> [stink].
\end{minted}
The variable \texttt{Num} is the same in \texttt{np(Num)} and \texttt{vp(Num)}, so the
query \texttt{phrase(s, [the, wumpus, stinks])} succeeds while \texttt{phrase(s, [the,
wumpus, stink])} fails. The query \texttt{phrase(s, X)} enumerates all (two) sentences of
the language. DCGs can be run online in SWISH, \url{https://swish.swi-prolog.org/}.
\end{nbexample}

\begin{nbintuition}[why arguments add power]
Without arguments, enforcing agreement needs one copy of each rule per combination of
features ($NP_{sing}$, $NP_{plur}$, \ldots). With arguments, one rule says ``whatever the
number is, it must be the same here and there''. Arguments can carry any Prolog term,
including counters, so DCGs can describe even languages that are not context-free, such
as $\mathbf{a}^n\mathbf{b}^n\mathbf{c}^n$.
\end{nbintuition}

\subsection{Statistical POS Tagging: Hidden Markov Models}
\label{sec:hmm-pos-tagging}
\slidesrc{Syntax: sentence, HMM: 1--20}

The rule-based approach of Section~\ref{sec:grammars-parsing} needs a linguist to write
the grammar. The statistical approach instead \emph{learns} from a corpus already tagged
by hand how likely each tag is after another tag, and how likely each word is for each
tag, then picks the most probable tag sequence for a new sentence. The classic model is
the Hidden Markov Model. The slides (from a Microsoft course on edX) introduce it with a
non-linguistic example, which we follow; then we map it to POS tagging.

\subsubsection{The model}
\slidesrc{Syntax: sentence, HMM: 3}

\begin{nbdef}{Hidden Markov Model (HMM)}
An \textbf{HMM} models a process that moves between \textbf{states}:
\begin{itemize}[nosep]
  \item the states are \textbf{hidden}: we cannot observe them;
  \item at each step the current state emits an \textbf{output} (observation), which we
        do see.
\end{itemize}
Our goal: from the sequence of outputs, find the most likely \textbf{sequence of
states}. Formally an HMM has
\begin{itemize}[nosep]
  \item a set of states $Q$ and of possible observations $O$;
  \item \textbf{start probabilities} $\pi(s)$: probability that the first state is $s$;
  \item \textbf{transition probabilities} $a(s' \to s) = P(s_t = s \mid s_{t-1} = s')$;
  \item \textbf{emission probabilities} $b_s(o) = P(o_t = o \mid s_t = s)$.
\end{itemize}
\end{nbdef}

Two assumptions make it ``Markov'' and keep it tractable: the next state depends only on
the current state (not on the whole history), and each output depends only on the state
that emits it.

\subsubsection{The weather example}
\slidesrc{Syntax: sentence, HMM: 4--5}

We are interested in the weather of a small town during the week of 5 October 1900.
No weather records were kept, but we have the diary of a teenager who wrote down the one
snack he had each day: hot chocolate, soda or ice cream. For our purposes the weather is
\emph{cold}, \emph{normal} or \emph{hot}. We are given:

\begin{center}
\small
\begin{tabular}{@{}l ccc @{\hspace{1.2cm}} l ccc@{}}
\toprule
\multicolumn{4}{@{}l}{\textbf{Start / transition} $P(\text{today} \mid \text{yesterday})$} &
\multicolumn{4}{l@{}}{\textbf{Emission} $P(\text{snack} \mid \text{weather})$}\\
\midrule
from $\downarrow$ to $\rightarrow$ & Cold & Normal & Hot & & hot choc. & soda & ice cream \\
\midrule
Start  & .1 & .8 & .1 & Cold   & .7 & .3 & 0 \\
Cold   & .7 & .1 & .2 & Normal & .1 & .7 & .2 \\
Normal & .3 & .4 & .3 & Hot    & 0  & .6 & .4 \\
Hot    & .2 & .4 & .4 & & & & \\
\bottomrule
\end{tabular}
\end{center}

Each row sums to 1. For example $P(\text{hot} \mid \text{cold}) = .2$ and
$P(\text{soda} \mid \text{hot}) = .6$. The diary for the 7 days is

\begin{center}
soda, soda, hot chocolate, ice cream, soda, soda, ice cream.
\end{center}

\textbf{Question}: what was the most likely weather sequence for the week?

\begin{nbintuition}[why not try every sequence?]
With 3 states and 7 days there are $3^7 = 2187$ possible weather sequences; with the 45
tags of a POS tagset and a 20-word sentence there are $45^{20}$, far too many. The
Viterbi algorithm finds the best one in time proportional to
(days) $\times$ (states)$^2$, using the same trick as the edit distance
(Section~\ref{sec:minimum-edit-distance}): for every state at every step, keep only the
best way to get there.
\end{nbintuition}

\subsubsection{The Viterbi algorithm}
\slidesrc{Syntax: sentence, HMM: 6--16}

\begin{nbfact}{Viterbi}
Let $v_t(s)$ be the probability of the \emph{best} state sequence that ends in state $s$
on day $t$ and emits the first $t$ observations.
\begin{align*}
  \text{day 1:}\quad v_1(s) &= \pi(s) \cdot b_s(o_1)
     \reason{p(transition) $\times$ p(output)}\\
  \text{day } t > 1\text{:}\quad v_t(s) &= \max_{s'} \bigl[\, v_{t-1}(s') \cdot a(s' \to s) \,\bigr] \cdot b_s(o_t)
\end{align*}
For each cell we also store a \textbf{backpointer} to the $s'$ that achieved the maximum.
At the end, pick the state with the largest $v_T(s)$ and follow the backpointers back to
day 1.
\end{nbfact}

\begin{nbexample}{the first two days}
\textbf{Day 1} (soda):
\begin{align*}
  v_1(\text{Cold})   &= .1 \cdot .3 = .03, &
  v_1(\text{Normal}) &= .8 \cdot .7 = .56, &
  v_1(\text{Hot})    &= .1 \cdot .6 = .06.
\end{align*}
\textbf{Day 2} (soda), cell Cold: which yesterday is the best way to arrive at Cold?
\begin{align*}
  v_2(\text{Cold}) &= \max\bigl(\underbrace{.03 \cdot .7}_{\text{from Cold}},\;
                                 \underbrace{.56 \cdot .3}_{\text{from Normal}},\;
                                 \underbrace{.06 \cdot .2}_{\text{from Hot}}\bigr) \cdot .3 \\
                   &= \max(.021,\ .168,\ .012) \cdot .3 = .168 \cdot .3 = .0504
     \reason{backpointer: Normal}
\end{align*}
Similarly $v_2(\text{Normal}) = \max(.03\cdot .1,\ .56\cdot .4,\ .06\cdot .4)\cdot .7 =
.224 \cdot .7 = .1568$ and $v_2(\text{Hot}) = \max(.03\cdot .2,\ .56\cdot .3,\ .06\cdot
.4)\cdot .6 = .168\cdot .6 = .1008$, both with backpointer Normal.
\end{nbexample}

Continuing the same way for all seven days fills the \textbf{trellis}:

\begin{center}
\footnotesize
\setlength{\tabcolsep}{4pt}
\begin{tabular}{@{}l ccccccc@{}}
\toprule
 & Day 1 & Day 2 & Day 3 & Day 4 & Day 5 & Day 6 & Day 7 \\
 & soda & soda & hot choc. & ice cream & soda & soda & ice cream \\
\midrule
Hot    & .06 & .1008 & 0 & \cellcolor{nbteal!25}\textbf{.00263} & .000632 & .000152 & \cellcolor{nbteal!25}\textbf{.0000248} \\
       & & {\scriptsize$\leftarrow$N} & & {\scriptsize$\leftarrow$C} & {\scriptsize$\leftarrow$H} & {\scriptsize$\leftarrow$H} & {\scriptsize$\leftarrow$N} \\
Normal & \cellcolor{nbteal!25}\textbf{.56} & \cellcolor{nbteal!25}\textbf{.1568} & .00627 & .000659 & \cellcolor{nbteal!25}\textbf{.000738} & \cellcolor{nbteal!25}\textbf{.000207} & .0000165 \\
       & & {\scriptsize$\leftarrow$N} & {\scriptsize$\leftarrow$N} & {\scriptsize$\leftarrow$C} & {\scriptsize$\leftarrow$H} & {\scriptsize$\leftarrow$N} & {\scriptsize$\leftarrow$N} \\
Cold   & .03 & .0504 & \cellcolor{nbteal!25}\textbf{.0329} & 0 & .000158 & .0000664 & 0 \\
       & & {\scriptsize$\leftarrow$N} & {\scriptsize$\leftarrow$N} & & {\scriptsize$\leftarrow$H} & {\scriptsize$\leftarrow$N} & \\
\bottomrule
\end{tabular}
\captionof{table}{The Viterbi trellis (values rounded to 3 significant digits). Under
each value, the backpointer: the state of the previous day that gave the maximum.
Shaded: the best path, recovered by backtracking from the largest value of day 7.}
\end{center}

\begin{itemize}[nosep]
  \item Some cells are 0 because the output is impossible in that state: no hot
        chocolate on a hot day, no ice cream on a cold day.
  \item On day 7 the best value is Hot (.0000248 $>$ .0000165), so the path ends in Hot.
        Its backpointer is Normal (day 6), whose backpointer is Normal (day 5), then Hot
        (day 4), then Cold (day 3), then Normal, Normal.
\end{itemize}

\begin{nbfact}{Result}
\textbf{Most likely weather}: normal, normal, cold, hot, normal, normal, hot.
\end{nbfact}

\begin{nbwarning}[rounding on the slides]
The slides compute with rounded intermediate values, so some cells differ slightly from
the exact ones (e.g.\ day 4 Hot: $.007 \cdot .4 = .0027$ instead of $.0329\cdot .2\cdot .4
= .00263$). One cell is a slip: day 2 Normal is written ``$.224 \cdot .7 = .1589$'', but
$.224\cdot .7 = .1568$. None of this changes the final path.
\end{nbwarning}

\begin{nbintuition}[numbers get tiny]
Every step multiplies by numbers smaller than 1, so after a few dozen words the values
underflow the floating-point range. Real implementations add \emph{logarithms} of the
probabilities instead of multiplying them: $\log$ is increasing, so the $\max$ (and the
best path) does not change.
\end{nbintuition}

\subsubsection{From weather to POS tags}
\slidesrc{Syntax: sentence, HMM: 17--20}

The weather problem is POS tagging in disguise:

\begin{center}
\small
\begin{tabular}{@{}l l l@{}}
\toprule
 & \textbf{Weather} & \textbf{POS tagging} \\
\midrule
hidden states    & cold, normal, hot & tags: DT, NN, VB, \ldots \\
observations     & snacks in the diary & words of the sentence \\
transitions      & $P(\text{hot} \mid \text{cold})$ & $P(\text{NN} \mid \text{DT})$: a noun after a determiner \\
emissions        & $P(\text{soda} \mid \text{hot})$ & $P(\text{flies} \mid \text{VBZ})$ \\
answer           & the weather of the week & the tag of each word \\
\bottomrule
\end{tabular}
\end{center}

Transition and emission probabilities are estimated by \emph{counting} in a hand-tagged
corpus (e.g.\ the Brown corpus): $P(\text{NN}\mid\text{DT}) = $ (number of times DT is
followed by NN) / (number of DTs). The Python code of the course implements Viterbi on
the weather example and prints the guessed sequence (\texttt{day=1 normal},
\texttt{day=2 normal}, \texttt{day=3 cold}, \ldots, \texttt{day=7 hot}).

\paragraph{HMM lab.} Experiment with the HMM Python code; extend it with other parts of
speech (see the Brown corpus tags,
\url{https://en.wikipedia.org/wiki/Brown_Corpus#Part-of-speech_tags_used}), words and
transitions, and use it to tag ``\emph{time flies like an arrow}''
(Section~\ref{sec:pos-tagging}); invent other short ambiguous sentences and extend the
code so that it tags them properly.

\subsubsection{Maximum Entropy Markov Models}
\slidesrc{Syntax: sentence, HMM: 19}

An HMM can reach very high accuracy in POS tagging, but it needs architectural
additions to deal with \textbf{unknown words}, suffixes and so on: an HMM can only use
$P(\text{word}\mid\text{tag})$, so a word never seen in training has no emission
probability at all.

\textbf{Maximum Entropy Markov Models} (MEMM) combine the HMM idea (a sequence of
hidden states, decoded with Viterbi) with \textbf{logistic regression}, a probabilistic
classifier trained with supervised machine learning. Instead of $P(\text{word} \mid
\text{tag})$, an MEMM directly models $P(\text{tag} \mid \text{word, previous tag,
features})$, where the features can be anything: suffix (\emph{-ing}, \emph{-ly}),
capitalisation, presence of digits or hyphens. This handles unknown words naturally. An
MEMM is trained on a POS-annotated text of the language; this is how the Stanford POS
tagger works (\url{https://nlp.stanford.edu/software/tagger.shtml}).

\subsubsection{Exercises}
\notinslides

\begin{nbexercise}{1}{a shorter diary}
With the same model, find the most likely weather for a two-day diary: \emph{hot
chocolate, ice cream}.
\end{nbexercise}
\begin{nbsolution}
Day 1 (hot chocolate): $v_1(\text{C}) = .1\cdot .7 = .07$, $v_1(\text{N}) = .8\cdot .1 =
.08$, $v_1(\text{H}) = .1 \cdot 0 = 0$.

Day 2 (ice cream):
\begin{align*}
  v_2(\text{C}) &= \max(\dots)\cdot 0 = 0\\
  v_2(\text{N}) &= \max(.07\cdot .1,\ .08\cdot .4,\ 0)\cdot .2 = .032\cdot .2 = .0064
     \reason{from N}\\
  v_2(\text{H}) &= \max(.07\cdot .2,\ .08\cdot .3,\ 0)\cdot .4 = .024\cdot .4 = .0096
     \reason{from N}
\end{align*}
Best final state: Hot (.0096), backpointer Normal. Answer: \textbf{normal, hot}.

\emph{Lesson}: on day 1 Cold explains hot chocolate better ($.7$ vs $.1$), but the start
probabilities favour Normal so much ($.8$ vs $.1$) that Normal wins. The best path is a
compromise between transitions and emissions.
\end{nbsolution}

\begin{nbsummary}[syntax at the sentence level]
\begin{itemize}[nosep]
  \item Sentence breaking: period = end, unless the previous token is an abbreviation,
        but yes if the next token is capitalised: about 95\% correct.
  \item POS tagging: tag each word using definition \emph{and} context; Penn Treebank
        tagset (36 tags: NN, NNS, VB, VBZ, JJ, DT, IN, \ldots).
  \item Generative grammar (Chomsky 1957). Chomsky hierarchy: regular $\subset$
        context-free $\subset$ context-sensitive $\subset$ recursively enumerable
        (finite automaton, pushdown automaton, linear-bounded TM, TM).
  \item CFG = lexicon + rules with one nonterminal on the left; parse tree; plain CFGs
        overgenerate (no agreement). DCGs add arguments to nonterminals (Prolog).
  \item HMM: hidden states, observed outputs; $\pi$, transitions $a$, emissions $b$.
        Viterbi: $v_t(s) = \max_{s'} v_{t-1}(s')\,a(s'\to s)\,b_s(o_t)$ + backpointers.
  \item MEMM = HMM + logistic regression; uses arbitrary features, handles unknown words.
\end{itemize}
\end{nbsummary}

